\documentclass[11pt,reqno]{amsart}

\usepackage{amsmath,amssymb,amsthm,mathtools}
\usepackage[T1]{fontenc}
\usepackage[utf8]{inputenc}
\usepackage{lmodern}
\usepackage{geometry}
\geometry{margin=2.8cm}
\usepackage{microtype}
\usepackage{hyperref}
\usepackage{cleveref}

\newtheorem{theorem}{Satz}[section]
\newtheorem{lemma}[theorem]{Lemma}
\newtheorem{corollary}[theorem]{Korollar}
\theoremstyle{definition}
\newtheorem{definition}[theorem]{Definition}
\theoremstyle{remark}
\newtheorem{remark}[theorem]{Bemerkung}

\numberwithin{equation}{section}

\title{Ein Titel für den mathematischen Aufsatz}
\author{Vorname Nachname}
\address{Institut für Mathematik, Beispiel-Universität}
\email{autor@example.org}
\date{\today}

\begin{document}

\begin{abstract}
  Hier steht die Zusammenfassung. Abstracts in \textsf{amsart} sind
  standardmäßig gesetzt und werden automatisch nummeriert.
\end{abstract}

\maketitle

\section{Einleitung}
\label{sec:einleitung}

Sei $X$ ein kompakter metrischer Raum. Wir bezeichnen mit
$\mathcal{C}(X)$ die Algebra der stetigen Abbildungen $X \to \mathbb{R}$
mit der Supremumsnorm.

\section{Vorbereitungen}
\label{sec:vorarbeit}

\begin{definition}[Norm]
  \label{def:norm}
  Eine Abbildung $\|\cdot\| : V \to \mathbb{R}$ heißt \emph{Norm}, wenn für alle
  $x,y \in V$ und $\lambda \in \mathbb{K}$ gilt:
  \begin{enumerate}
    \item $\|x\| \ge 0$ mit $\|x\| = 0 \Leftrightarrow x = 0$,
    \item $\|\lambda x\| = |\lambda|\,\|x\|$,
    \item $\|x + y\| \le \|x\| + \|y\|$.
  \end{enumerate}
\end{definition}

\begin{theorem}[Banach-Fixpunktsatz]
  \label{thm:banach}
  Sei $(X,d)$ ein vollständiger metrischer Raum und $T : X \to X$ eine
  Kontraktion mit Konstante $q \in [0,1)$. Dann besitzt $T$ genau einen
  Fixpunkt $x^\ast$ und für jedes $x_0 \in X$ konvergiert die Folge
  $(x_n)$ mit $x_{n+1} = T(x_n)$ gegen $x^\ast$ mit
  \begin{equation}
    \label{eq:konvergenz}
    d(x_n, x^\ast) \le \frac{q^n}{1-q}\, d(x_1, x_0).
  \end{equation}
\end{theorem}

\begin{proof}
  Nach \cref{def:norm} genügt es, die Folge $(x_n)$ als Cauchy-Folge zu
  zeigen. Wir schätzen für $m > n$:
  \begin{align*}
    d(x_n, x_m) &\le \sum_{k=n}^{m-1} d(x_k, x_{k+1})
                 \le \sum_{k=n}^{m-1} q^k\, d(x_0, x_1) \\
                &\le \frac{q^n}{1-q}\, d(x_0, x_1) \xrightarrow[n\to\infty]{} 0.
  \end{align*}
  Die Stetigkeit von $T$ liefert $T(x^\ast) = x^\ast$. \qed
\end{proof}

\begin{remark}
  Die Abschätzung \cref{eq:konvergenz} ist schneidengenau~\cite{rudin1976}.
\end{remark}

\begin{corollary}
  \label{cor:loesung}
  Die Gleichung $x = T(x)$ besitzt in $\mathbb{R}^n$ eine eindeutige Lösung,
  falls $\|DT(x)\|_\infty \le q < 1$ auf einer abgeschlossenen Kugel gilt.
\end{corollary}

\begin{thebibliography}{9}
  \bibitem{rudin1976} W. Rudin, \emph{Principles of Mathematical Analysis},
  3rd ed., McGraw-Hill, 1976.
\end{thebibliography}

\end{document}
